% 原书第六篇6.6节259--266；
% 印刷261--262，PDF270--271，书末答案350（PDF359）。
\nuclearreactionproblem{259}{261}
用原子质量单位计算：
(a) $\mathrm{^{8}Li}$ 原子质量，其核结合能
$E_b=41.3\,\mathrm{MeV}$；
(b) $\mathrm{^{10}C}$ 原子核质量，其平均每核子
结合能为 $6.04\,\mathrm{MeV}$。
\exsolution
取 $1u\,c^2=931.5\,\mathrm{MeV}$，
原书氢原子质量 $m_H=(1+0.00783)u$、
中子质量 $m_n=(1+0.00867)u$。
(a) Li 有 $Z=3,A=8$，
\[
 \boxed{M_{\mathrm{^{8}Li\,atom}}=
 3m_H+5m_n-\frac{41.3\,\mathrm{MeV}}{c^2}
 \simeq8.0225u}.
\]
(b) C 有 $Z=6,A=10$，
结合能 $E_b=10(6.04)=60.4\,\mathrm{MeV}$。
求的是\emph{原子核}而不是原子，
所以应使用裸质子质量
$m_p=m_H-m_e$：
\[
 \boxed{M_{\mathrm{^{10}C\,nucleus}}=
 6(m_H-m_e)+4m_n-\frac{60.4\,\mathrm{MeV}}{c^2}
 \simeq10.0135u}.
\]
\excheck (b) 若直接代氢\emph{原子}质量而不扣去6个电子，
将错误地计算中性碳原子的质量。

\nuclearreactionproblem{260}{261}
某核反应
$A_1+A_2\to A_3+A_4$，
四核的结合能分别为 $E_1,E_2,E_3,E_4$。
求反应能 $Q$。
\exsolution
核反应能量守恒
$Q=(M_1+M_2-M_3-M_4)c^2$。
各核质量由
$M_ic^2=Z_im_pc^2+N_im_nc^2-E_i$。
在电荷数与核子数守恒下，
所有自由核子的总静能在反应前后抵消，
因此
\[
 \boxed{Q=(E_3+E_4)-(E_1+E_2)}.
\]
\excheck 产物的总结合能大于反应物时，
$Q>0$，说明放能；反之吸能。

\nuclearreactionproblem{261}{261--262}
假设每个 $\mathrm{^{235}U}$ 核裂变释放
$200\,\mathrm{MeV}$。
(a) 求1kg $\mathrm{^{235}U}$ 裂变释放的能量，
以及等能量、热值$30\,\mathrm{kJ/g}$的煤质量；
(b) 30千吨TNT当量的核爆炸对应多少裂变铀质量？
TNT热值为$4.1\,\mathrm{kJ/g}$。
\exsolution
(a) 1kg铀的核数
$N=(1000\,\mathrm g/235\,\mathrm{g/mol})N_A$，
释放总能量
\[
 \boxed{E=N(200\,\mathrm{MeV})
 \simeq8.21\times10^{13}\,\mathrm J
 =8.21\times10^{10}\,\mathrm{kJ}}.
\]
煤热值
$30\,\mathrm{kJ/g}=3.0\times10^7\,\mathrm{J/kg}$，
故
\[
 \boxed{m_{\rm coal}=
 \frac E{3.0\times10^7\,\mathrm{J/kg}}
 \simeq2.74\times10^6\,\mathrm{kg}}.
\]
(b) 一千吨TNT当量等于 $10^6\,\mathrm{kg}$ TNT，
30千吨即 $3.0\times10^7\,\mathrm{kg}$。
其能量
$E_{\rm TNT}=3.0\times10^7(4.1\times10^6)\,\mathrm J
=1.23\times10^{14}\,\mathrm J$。
因此等能铀质量
\[
 \boxed{m_U=\frac{E_{\rm TNT}}
 {8.21\times10^{13}\,\mathrm{J/kg}}
 \simeq1.50\,\mathrm{kg}}.
\]
\excheck 本版书末(a)煤质量写
$2.7\times10^8\,\mathrm{kg}$，
但依原题 $30\,\mathrm{kJ/g}$
计算应为$2.7\times10^6\,\mathrm{kg}$，
差100倍。两结果保留校勘记录。

\nuclearreactionproblem{262}{262}
由氘核聚变生成1.0g $\mathrm{^{4}He}$
时释放多少热？该热量相当于热值
$30\,\mathrm{kJ/g}$ 的多少煤？
\exsolution
反应
$2\,\mathrm{^{2}H}\to\mathrm{^{4}He}+\gamma$
（这里只比较核静质量）。
原书原子质量超额
$\Delta_D=0.01410u$、
$\Delta_\alpha=0.00260u$，
因此
\[
 Q=(2\Delta_D-\Delta_\alpha)c^2
 =0.02560u\,c^2
 \simeq23.85\,\mathrm{MeV}.
\]
1克氦核数
$N=(1.0\,\mathrm g/4\,\mathrm{g/mol})N_A$，
故
\[
 \boxed{E=NQ\simeq5.75\times10^{11}\,\mathrm J
 =5.75\times10^8\,\mathrm{kJ}}.
\]
煤热值
$30\,\mathrm{kJ/g}=3.0\times10^7\,\mathrm{J/kg}$，
\[
 \boxed{m_{\rm coal}\simeq1.92\times10^4\,\mathrm{kg}}.
\]
\excheck 原书书末写
$5.74\times10^7\,\mathrm{kJ}$，
但同一行给煤约$2\times10^4\,\mathrm{kg}$；
根据质量亏损及原题热值，能量应为
$5.74\times10^8\,\mathrm{kJ}$，书末相差10倍。

\nuclearreactionproblem{263}{262}
利用原子质量表，
求反应
$\mathrm{^{6}Li}+\mathrm{^{2}H}\to2\alpha$
每个核子释放的能量，
并与 $\mathrm{^{235}U}$ 每次裂变放200MeV时的
每核子放能量比较。
\exsolution
原书质量超额（单位$u$）
$\Delta_{\mathrm{^{6}Li}}=0.01513$、
$\Delta_D=0.01410$、
$\Delta_\alpha=0.00260$。
故
\[
 Q=[0.01513+0.01410-2(0.00260)]
 (931.5\,\mathrm{MeV})
 \simeq22.38\,\mathrm{MeV}.
\]
反应参与核子总数8，
所以每核子放能
\[
 \boxed{Q/8\simeq2.80\,\mathrm{MeV/nucleon}}.
\]
铀裂变每核子平均放能
$\boxed{200/235\simeq0.851\,\mathrm{MeV/nucleon}}$，
二者比约3.3。
\excheck 这里按\emph{参与反应的全部8个核子}
而不是按两个产物的\emph{个数}来归一化。

\nuclearreactionproblem{264}{262}
$\mathrm{^{7}Li}+p\to2\alpha$，
已知 $\mathrm{^{7}Li}$、
$\mathrm{^{4}He}$ 每核子结合能分别为
$5.60$ 与 $7.06\,\mathrm{MeV}$。
求核反应能。
\exsolution
反应初态总结合能
$E_i=7(5.60)+0$，自由质子结合能为零。
末态
$E_f=2\cdot4(7.06)$。
反应能是结合能差
\[
 \boxed{Q=E_f-E_i=
 8(7.06)-7(5.60)\simeq17.3\,\mathrm{MeV}}.
\]
\excheck 产物总结合能大于初态，
故该反应放能。

\nuclearreactionproblem{265}{262}
反应
$\mathrm{^{14}N}(\alpha,p)\mathrm{^{17}O}$
中入射 $\alpha$ 动能
$T_\alpha=4.0\,\mathrm{MeV}$，
出射质子动能
$T_p=2.09\,\mathrm{MeV}$，
质子与入射方向夹角
$\theta=60^\circ$。
求反应能 $Q$。
\exsolution
以实验室系静靶核为参考，
动量守恒
$\boldsymbol p_\alpha=\boldsymbol p_p+\boldsymbol p_O$，
所以
$p_O^2=p_\alpha^2+p_p^2-2p_\alpha p_p\cos\theta$。
非相对论动量
$p^2=2mT$，故反冲氧核动能
\[
 T_O=\frac{m_\alpha T_\alpha+m_p T_p
 -2\sqrt{m_\alpha m_pT_\alpha T_p}\cos\theta}{m_O}.
\]
近似质量数 $m_\alpha:m_p:m_O=4:1:17$，
得
$T_O\simeq0.724\,\mathrm{MeV}$。
根据动能能量守恒
$Q=T_p+T_O-T_\alpha$，
\[
 \boxed{Q\simeq2.09+0.724-4.00
 \simeq-1.19\,\mathrm{MeV}\approx-1.2\,\mathrm{MeV}}.
\]
\excheck 质子出射角相对\emph{入射束}，
故反冲项是
$-2p_\alpha p_p\cos\theta$；
该反应为吸能反应。

\nuclearreactionproblem{266}{262}
利用原子质量表计算：
(a)$\mathrm{^{7}Li}(p,n)\mathrm{^{7}Be}$；
(b)$\mathrm{^{9}Be}(n,\gamma)\mathrm{^{10}Be}$；
(c)$\mathrm{^{7}Li}(\alpha,n)\mathrm{^{10}B}$；
(d)$\mathrm{^{16}O}(d,\alpha)\mathrm{^{14}N}$
四个核反应的反应能。
\exsolution
对守恒总核子数的反应，
用\emph{中性原子质量超额}
求静质量差
$Q=(\sum\Delta_{\rm in}-\sum\Delta_{\rm out})c^2$，
其中
$1u\,c^2=931.5\,\mathrm{MeV}$。
从原书表取
$\Delta_H=0.00783$、
$\Delta_n=0.00867$、
$\Delta_{\mathrm{^{7}Li}}=0.01601$、
$\Delta_{\mathrm{^{7}Be}}=0.01693$、
$\Delta_{\mathrm{^{9}Be}}=0.01219$、
$\Delta_{\mathrm{^{10}Be}}=0.01354$、
$\Delta_{\mathrm{^{4}He}}=0.00260$、
$\Delta_{\mathrm{^{10}B}}=0.01294$、
$\Delta_{\mathrm{^{16}O}}=-0.00509$、
$\Delta_D=0.01410$、
$\Delta_{\mathrm{^{14}N}}=0.00307$，单位均为$u$。
于是
\[
 \boxed{Q_a=(0.01601+0.00783-0.00867-0.01693)c^2
 \simeq-1.64\,\mathrm{MeV}},
\]
\[
 \boxed{Q_b=(0.01219+0.00867-0.01354)c^2
 \simeq+6.82\,\mathrm{MeV}},
\]
\[
 \boxed{Q_c=(0.01601+0.00260-0.00867-0.01294)c^2
 \simeq-2.79\,\mathrm{MeV}},
\]
\[
 \boxed{Q_d=(-0.00509+0.01410-0.00260-0.00307)c^2
 \simeq+3.11\,\mathrm{MeV}}.
\]
\excheck $\gamma$ 静质量为零；
全部反应电荷守恒，
可以用同一中性原子质量超额记法互相抵消电子质量。
